\section{From formal sign fields to a smooth splitting}
\label{smooth:section}

The finite alternative in Theorem~\ref{alg:dichotomy} remains.  In this
section all work on the universal cover is local in the resulting smooth
structure, so the conclusions descend to $SM$.  We first recover the actual
plaque jets on an entire typical axis.  We then propagate their regularity
using scalar functions which label weak unstable leaves.

\subsection{Actual plaque jets and formal leaf labels}

\begin{lemma}[Tangency to the central plaques]
\label{smooth:tangency}
In the finite alternative, the positive-$D$ invariant formal plane field is
tangent to the entire formal central unstable plaque.  Consequently, on each
entire typical strong stable leaf, all finite Taylor jets of the actual weak
unstable plaques are smooth functions of their crossing point.  They have
holonomic full formal extensions there.  The same statements hold on the
weak stable leaf generated by that strong stable leaf.
\end{lemma}

\begin{proof}
Use the graded section coordinates of Lemma~\ref{sym:graded-contact}
and the dilation of Proposition~\ref{sym:dilation}.  Write $x$ for the positive variables and $y$ for
the negative variables.  The actual central unstable plaque has formal
projection $y=0$.  Let $D_+$ and $D_-$ denote the restrictions of $D$ to the
two coordinate spaces.  The positive formal field furnished by
Proposition~\ref{alg:axis-coherence} has value $E^u$ at the origin and is invariant
under the dilation.  Its restriction to $y=0$ is a graph, with graph matrix
$B(x)$ from the positive to the negative space.  Dilation invariance gives
\[
 B(e^{tD_+}x)=e^{tD_-}B(x)e^{-tD_+}.
\]
Diagonalize $D$ over the real graded spaces.  A monomial $x^\nu$ in a matrix
entry with input weight $\lambda_i>0$ and output weight $\mu_j<0$ would
have to satisfy
\[
 \sum_a\nu_a\lambda_a=\mu_j-\lambda_i<0.
\]
The left side is nonnegative, so every coefficient vanishes.  Thus $B=0$.
This is an identity of formal series; it needs no convergence.  Adding the
flow direction gives formal tangency to the weak unstable plaque.

In ordinary smooth coordinates, represent that plaque by a graph $z=g(v)$
with prescribed crossing $g(0)$.  Tangency to the formal plane field gives a
system
\[
 \partial_{v_i}g(v)=A_i(v,g(v)).
\]
Its constant coefficient determines the first derivatives of $g$; after
$d$ differentiations, it determines derivatives of order $d+1$ from the
field jet and the already determined graph coefficients.  There is no
existence assertion hidden in this recursion: the actual central plaque
already provides the formally tangent solution.  The recursion merely
expresses each of its graph coefficients by finitely many operations on
the ambient formal field jet.

The full ambient jets of that field are smooth and holonomic on each
entire typical stable axis, by Proposition~\ref{alg:axis-coherence}.  Hence the
recursively obtained graph coefficients are smooth there and have the
same tangential coherence.  Initially the equality with actual plaque
coefficients holds at the regular crossing points.  It holds everywhere
on the axis by continuity: the actual plaques depend continuously in
$C^0$ on their crossings, have uniformly bounded local derivatives of
each finite order, and therefore depend continuously in each finite
$C^r$ norm by interior interpolation.  Regular crossings are dense on the
axis by conditional full measure.  This also explains why no transverse
smoothness of the actual foliation was used.  Flow transport proves the
weak-axis assertion, with compatibility on overlapping flow boxes.
\end{proof}

Fix one such strong stable leaf $A^{\mathrm{all}}$ on the universal cover,
with future endpoint $\eta$.  Normalize its Busemann function so that
$A^{\mathrm{all}}$ is at height $b_\eta=0$; our convention is
$b_\eta(\pi\phi_t v)=b_\eta(\pi v)-t$ when $v_+=\eta$.
The past-endpoint map identifies this leaf with
$\partial_\infty\widetilde M\setminus\{\eta\}$.  Choose a coordinate patch
$A\subset A^{\mathrm{all}}$, with smooth coordinates $a\mapsto q_A(a)\in
\R^m$, and write
\[
 \mathcal B=\{\phi_t a:a\in A,\ t\in\R\}.
\]
This is an open part of the fixed-future weak stable leaf.  On the open
set $\mathcal U$ of vectors whose past endpoints occur in $A$, define
\begin{equation}
 \label{smooth:label-definition}
 q(v)=q_A(a(v)),\qquad a(v)_-=v_-,\quad a(v)\in A.
\end{equation}
The endpoint identifications show that $q$ is continuous.  It is constant
on actual weak unstable leaves and under the flow.  Only its restriction
$q_0=q|_{\mathcal B}$ is known to be smooth at this stage.

\begin{lemma}[Formal inversion of plaque labels]
\label{smooth:formal-labels}
The function $q_0$ has a distinguished full formal ambient extension along
$\mathcal B$.  Its coefficients are smooth and holonomic on $\mathcal B$,
it is formally constant on every actual weak unstable plaque crossing
$\mathcal B$, and its finite jets transform by scalar pullback under the
flow.  These assertions concern formal jets along $\mathcal B$, and do
not assert differentiability of $q$ elsewhere.
\end{lemma}

\begin{proof}
At any fixed Busemann height, the strong stable slice in $\mathcal B$ is a
smooth transversal to the weak unstable plaques.  Use coordinates $a$ on
this slice and auxiliary variables $z$ along the plaques.  By
Lemma~\ref{smooth:tangency}, the parameterized actual plaque Taylor graphs
form a formal map
\[
 \Psi(a,z),\qquad \Psi(a,0)=a,
\]
whose coefficients in $z$ are smooth in $a$, with all their tangential
derivatives. Within a weak unstable plaque, the strong plaque through
its crossing is tangent to $\ker\alpha$. Its graph coefficients are
determined recursively from the weak plaque jets and the smooth contact
form, as in Lemma~\ref{smooth:tangency}; the actual strong plaque supplies
the solution. These finite operations preserve smooth dependence on $a$.
We may therefore choose $z=(u,t)$ and write
\[
 \Psi(a,u,t)=\phi_t\Gamma(a,u),\qquad \Gamma(a,0)=a,
\]
where $\Gamma(a,u)$ is the actual strong unstable plaque Taylor map
through $a$. Only its finite Taylor coefficients and their smooth
dependence on $a$ are used. The differential of $(a,z)\mapsto\Psi(a,z)$ at $z=0$ is
invertible, by transversality.  Formal inversion therefore determines
$a=a(p)$ and the scalar formal function $q_0|_{\mathrm{slice}}(a(p))$
to every finite order.  Concretely, to compute order $r$, take a sufficiently high finite
Taylor polynomial of $\Psi$ in $z$ and invert it as an ordinary local
smooth map.  The resulting order-$r$ jet is independent of all higher
coefficients.  Thus this construction uses only finitely many jet
operations at each stage.

Smooth dependence of the coefficients on $a$ gives tangential
holonomicity of the inverse jets.  The construction is independent of
coordinates because it describes the crossing of an actual plaque with
the chosen transversal.  Moving the height slice by the flow transports
both the actual plaque jets and the prescribed labels; formal inversion
commutes with this transport.  This proves compatibility at different
heights. The inverse label is independent of $u,t$, and
$\partial_t\Psi=X\circ\Psi$ gives the full formal identity $Xq=0$.
Differentiating scalar jet equivariance in time therefore gives the
holonomic derivative identity along $X$; combined with coherence on
each height slice, this proves full holonomicity along $\mathcal B$.
In particular, the scalar jet identity is
\[
 j^r_p q=(j^r_{\phi_t p}q)\circ j^r_p\phi_t,
 \qquad p,\phi_t p\in\mathcal B,
\]
where both $q$-jets in this display mean the formal jets just constructed.
\end{proof}

\subsection{Backward geometry and estimates along the axis}

Let $N_\eta(x)=(x,v_\eta(x))$ be the unit vector based at $x$ pointing
toward $\eta$.  It parametrizes the fixed-future weak stable leaf as a
smooth graph over $\widetilde M$.  This graph, its local charts, and all
of its fixed-order local derivatives have uniform bounds.  Indeed the
individual weak stable leaves have uniform smooth plaque bounds, and
their spatial projections have uniformly invertible differentials.
Write $\rho(x,v)=N_\eta(x)$.  Near this graph use the intrinsic angular
coordinate
\[
 t= v-\langle v,v_\eta(x)\rangle v_\eta(x)
       \in v_\eta(x)^\perp,
 \qquad \langle v,v_\eta(x)\rangle>0.
\]
It is the tangential component of $v$ in the horosphere.  Thus
$v=\sqrt{1-|t|^2}\,v_\eta(x)+t$.  The degree-$l$ angular Taylor
coefficient of the formal label will be denoted $q_l$; it is a smooth
symmetric $l$-tensor on the horosphere tangent bundle over $\mathcal B$,
with values in $\R^m$, and includes the Taylor factorial.  In these
coordinates the formal label is $\sum_{l\geq0}q_l(s,r)[t^l]$, where $r$
is Busemann height and $s$ denotes coordinates on its horosphere.

\begin{lemma}[Compact sets approach the fixed-future axis]
\label{smooth:backward-geometry}
Let $K\Subset\mathcal U$.  There are $b>0$, $C<\infty$, and a compact
$K_{\mathcal B}\Subset\mathcal B$ such that, for all sufficiently large
$n$ and $v\in K$,
\[
 |t(f^{-n}v)|\leq Ce^{-bn},\qquad
 f^n\rho(f^{-n}v)\in K_{\mathcal B}.
\]
The actual weak unstable plaque through $f^{-n}v$ crosses the axis
$t=0$ at the same Busemann height, at a point $a_n(v)$.  That crossing
and the axis points underlying the intervening local plaque graph also
have their $f^n$-images in $K_{\mathcal B}$, after enlarging this compact
set.  All assertions hold on a fixed neighborhood of $K$.
\end{lemma}

\begin{proof}
The crossing $a(v)$ in \eqref{smooth:label-definition} depends
continuously on $v$, so $a(K)$ is compact in $A$.  Put $\zeta=v_-$ and
choose the continuous bounded time shift
\[
 \tau(v)=b_\zeta(\pi v)-b_\zeta(\pi a(v)).
\]
The vector $z(v)=\phi_{\tau(v)}a(v)$ has the same past endpoint and the
same past Busemann height as $v$; hence it is on its strong unstable
leaf.  Moreover their intrinsic strong unstable distance is uniformly
bounded on $K$.  To see the uniformity explicitly, join their base
points by a spatial geodesic segment of bounded length and project it
onto their common past horosphere along the Busemann gradient flow.
The required flow times are bounded by the segment length.  Uniform
bounds for the Busemann Hessian bound the length of this projected path
and its lift to the normal graph.  This uses the fixed negatively
curved metric, not regularity of the endpoint map.

Backward strong unstable contraction now gives
\[
 d^u(f^{-n}v,f^{-n}z(v))\leq Ce^{-bn}.
\]
Write $f^{-n}v=(x_n,v_n)$ and $f^{-n}z(v)=N_\eta(y_n)$.
The uniform local bounds of $N_\eta$ imply
$d(N_\eta(x_n),N_\eta(y_n))\leq Ce^{-bn}$ and the asserted angular
bound.  If
\[
 \epsilon_n=b_\eta(x_n)-b_\eta(y_n),
\]
then $|\epsilon_n|\leq Ce^{-bn}$.  The two vectors
$\phi_{\epsilon_n}N_\eta(x_n)$ and $N_\eta(y_n)$ belong to the same
strong stable leaf and have intrinsic distance $O(e^{-bn})$.
Forward strong stable contraction, and the unchanged small time shift,
show that $f^nN_\eta(x_n)$ stays within $O(e^{-bn})$ in the intrinsic
weak stable metric of the compact family $z(v)$.  This proves the
compact-axis assertion.

For completeness, the local crossing assertion has a uniform
transversality bound.  At $t=0$, the weak unstable plaque intersected with
$b_\eta=r$ has tangent $E^u$.  In Jacobi coordinates the differential of
its angular coordinate in a spatial direction $a$ is
$(U(v)+U(-v))a$, which is uniformly positive definite.  Uniform smooth
bounds on individual plaques and the inverse function theorem therefore
give actual graphs
\begin{equation}
 \label{smooth:actual-graphs}
 s=P_{s_0,r}(t),\qquad P_{s_0,r}(0)=s_0,
\end{equation}
on a fixed radius, with uniformly bounded derivatives in $t$ of each
fixed order.  No derivative with respect to $s_0$ is asserted.

The nearby crossing of the plaque through $f^{-n}v$ is especially
explicit.  If $c_n=b_\eta(x_n)$, then
\[
 c_n=nT-\tau(v)+\epsilon_n,
 \qquad a_n(v)=\phi_{-c_n}a(v),
 \qquad f^na_n(v)=\phi_{\tau(v)-\epsilon_n}a(v).
\]
These identities follow from its past endpoint, future endpoint, and
height.  They prove compactness of its forward image.  On the local
graph between angular values $0$ and $t(f^{-n}v)$, the base points move
by $O(e^{-bn})$ on this same horosphere.  Their axis vectors have that
small intrinsic stable distance from $a_n(v)$; contraction gives the
remaining compactness assertion.  Replacing $K$ by a slightly larger
compact subset of $\mathcal U$ proves the neighborhood assertion.
\end{proof}

All following norms along the axis are computed in uniformly bounded
local weak stable charts and uniformly bounded tensor frames.  For
$L\Subset\mathcal B$ let $L_n=f^{-n}L$.  When an interior estimate is
used, enlarge $L$ slightly inside $\mathcal B$. The forward weak-plaque
bounds in Lemma~\ref{pre:plaque-contraction}, in particular the uniform
Lipschitz bound, give a uniform positive amount of local room around $L_n$.

\begin{lemma}[A derivative-order-independent exponential bound]
\label{smooth:naive}
For each angular degree $l$ there is $A_l<\infty$ such that, for every
$j\geq0$ and every $L\Subset\mathcal B$,
\begin{equation}
 \label{smooth:naive-bound}
 \|q_l\|_{C^j(L_n)}\leq C_{L,l,j}e^{A_ln}.
\end{equation}
The exponent $A_l$ is independent of $j$.
\end{lemma}

\begin{proof}
Transport the full ambient scalar jet of order $l$ from $f^np$ to $p$,
using Lemma~\ref{smooth:formal-labels}.  The one-step transport is a
finite-dimensional linear pullback-jet matrix $M_l(p)$.  In uniformly
bounded geometric charts these matrices and every fixed-order weak-axis
derivative are uniformly bounded: the one-step ambient map is smooth,
and the weak-axis charts have the uniform leafwise bounds just
explained.  The transport over $n$ steps is the product of these matrices
at $p,fp,\ldots,f^{n-1}p$.

When this product is differentiated $j$ times along the initial weak
axis, at most $j$ factors are differentiated.  Each derivative of a
factor composed with $f^i|_{\mathcal B}$ is bounded by a constant
depending on $l,j$, uniformly in $i$.  The remaining factors are bounded
by a fixed $M_l\geq1$.  The number of terms is at most a constant times
$(1+n)^j$.  The final jet field on $L$, and its derivatives composed with
$f^n|_{\mathcal B}$, have bounded derivatives.  Hence the product rule
gives $C_{L,l,j}(1+n)^jM_l^n$.  Absorb the polynomial in $e^n$, with its
$j$-dependent prefactor.  Taking angular coefficients and changing
uniform frames preserve the resulting exponent $A_l=1+\log M_l$.
This is Lemma~\ref{real:product-derivatives} on the fixed bundle
$J^l$ of scalar $l$-jets: differentiation changes the matrix
coefficient functions, not the jet order. Using ambient derivative bounds
for $f^n$ in its place would lose the independence from $j$.
\end{proof}

\begin{lemma}[Subexponential bounds for every axis derivative]
\label{smooth:subexponential}
For every $l,j\geq0$, every $\varepsilon>0$, and every
$L\Subset\mathcal B$ there is a constant $C_{L,l,j,\varepsilon}$ such
that
\begin{equation}
 \label{smooth:subexponential-bound}
 \|q_l\|_{C^j(L_n)}
 \leq C_{L,l,j,\varepsilon}e^{\varepsilon n}.
\end{equation}
\end{lemma}

\begin{proof}
We induct on $l$, proving all derivative orders at each stage.  For
$l=0$, the identity $q_0(p)=q_0(f^np)$ and the bounded weak-axis
derivatives of forward iterates give bounded derivatives on $L_n$.

At each fixed height, formal constancy on the actual graph
\eqref{smooth:actual-graphs} gives the coefficientwise identity
\begin{equation}
 \label{smooth:crossing-identity}
 \sum_{k\geq0}q_k(P_{s_0,r}(t),r)[t^k]=q_0(s_0,r).
\end{equation}
Every occurrence of a coefficient on the left is Taylor-expanded at
$t=0$.  The identity follows directly from the formal crossing inversion
of Lemma~\ref{smooth:formal-labels} and its tangential holonomicity.
In degree $l$, its only term containing $q_l$ is
$q_l(s_0,r)[t^l]$.  All other terms involve derivatives in $s$ of $q_k$
with $k<l$ and order at most $l-k$, multiplied by polynomials in the
$t$-derivatives of $P_{s_0,r}$ through order $l$.  These graph derivatives
are uniformly bounded.  The induction hypothesis and polarization of
the homogeneous degree-$l$ identity therefore give
\[
 \|q_l\|_{C^0(L_n)}\leq C_{L,l,\varepsilon}e^{\varepsilon n}
 \quad\hbox{for every }\varepsilon>0.
\]
This argument estimates only values of the next coefficient.  In
particular it never differentiates the actual graph with respect to its
crossing point.

We upgrade values to arbitrary weak-axis derivatives using interpolation
on the full weak axis, including the height coordinate.  On a fixed
interior chart the elementary Taylor difference estimate is
\[
 \|D^ju\|_{C^0}
 \leq C_{j,r}\bigl(h^{-j}\|u\|_{C^0}
                 +h^{r-j}\|u\|_{C^r}\bigr),
 \qquad 0<h<h_0,
\]
where the norms on the right are on a slightly larger chart.  The
uniform interior room makes its constants independent of $n$; see also
Lemma~\ref{real:interpolation}.  Fix $j\geq1$ and target
$\varepsilon>0$.  Take $\delta=\varepsilon/(4j)$ and
$h=e^{-\delta n}$, use the value bound with exponent $\varepsilon/4$,
and choose a fixed integer $r>j+A_l/\delta$.  By
\eqref{smooth:naive-bound}, the two terms then have exponents at most
$\varepsilon/2$ and $A_l-\delta(r-j)<0$, respectively.  This proves
\eqref{smooth:subexponential-bound}.  The choice of $r$ is possible
precisely because $A_l$ does not depend on $r$.  The argument works for
every compact $L$, using a larger compact for interior estimates, so the
induction is complete.
\end{proof}

\subsection{Smoothness of the actual label}

\begin{proposition}[Propagation from the axis]
\label{smooth:propagation}
The actual endpoint label $q:\mathcal U\to\R^m$ is smooth.
\end{proposition}

\begin{proof}
Fix nested compact neighborhoods $K\Subset K^+\Subset\mathcal U$.
For large $n$, Lemma~\ref{smooth:backward-geometry} puts $f^{-n}K^+$ in
the angular coordinate neighborhood of the axis.  For a fixed finite
cutoff $R$ define the smooth function, near $K^+$,
\begin{equation}
 \label{smooth:approximants}
 p_{n,R}(v)=\sum_{l=0}^R
 q_l\bigl(\rho(f^{-n}v)\bigr)
       \bigl[t(f^{-n}v)^l\bigr].
\end{equation}
These expressions are intrinsic tensor contractions, so no seams from
angular frames occur.  They use a finite polynomial at each $n$; no
infinite formal series is evaluated.

We first estimate their error without differentiating $q$.  At a pulled
back point, use its actual plaque graph
$s=P_{s_0,r}(t)$ and its genuine crossing $(s_0,r,0)$.  The function of
the angular parameter
\[
 H(t)=\sum_{l=0}^R q_l(P_{s_0,r}(t),r)[t^l]-q_0(s_0,r)
\]
has zero Taylor coefficients through degree $R$ by
\eqref{smooth:crossing-identity}.  Along the segment from zero to the
angular parameter in question, all underlying axis points lie over the
compact set supplied by Lemma~\ref{smooth:backward-geometry}.
Lemma~\ref{smooth:subexponential}, the uniform derivatives of the
single actual graph, and Taylor's remainder formula give
\[
 |H(t)|\leq C_{K,R,\varepsilon}e^{\varepsilon n}|t|^{R+1}.
\]
The actual label equals $q_0(s_0,r)$ on that plaque, and is flow
invariant.  Consequently
\begin{equation}
 \label{smooth:approximation-error}
 \|p_{n,R}-q\|_{C^0(K^+)}
 \leq C_{K,R,\varepsilon}
       e^{-(b(R+1)-\varepsilon)n}.
\end{equation}
Only derivatives of the smooth axis coefficients and of an individual
actual plaque entered this estimate.

For each fixed $j$, the same coefficients have subexponential
$C^j$ bounds.  A polynomial of degree at most $R$ in $|t|<1$ therefore
has a $C^j$ bound $C_{K,R,j,\varepsilon}e^{\varepsilon n}$ in the
geometric angular coordinates.  The ordinary ambient bounds for
$f^{-n}$ and the chain rule give
\begin{equation}
 \label{smooth:approximant-growth}
 \|p_{n,R}\|_{C^j(K^+)}
 \leq C_{K,R,j,\varepsilon}e^{(B_j+\varepsilon)n},
\end{equation}
where $B_j$ depends on $j$ and the smooth dynamics, but is independent
of $R$.  Indeed differentiating the angular polynomial at most $j$
times only introduces $R$-dependent combinatorial constants; its
coefficients already have arbitrary subexponential slack.  All
exponential losses come from derivatives of $f^{-n}$ of orders at most
$j$.  This independence from $R$ is the second essential exponent
separation in the proof.

Here are the convergence quantifiers.  Fix a desired derivative order
$j\geq1$ and a desired positive convergence rate $\lambda$.  Set
$r=2j+2$, take $\varepsilon=1$, and choose one finite $R$ large enough
that
\[
 \bigl(b(R+1)-1\bigr)(1-j/r)
       -(B_r+1)j/r>\lambda.
\]
The differences $p_{n+1,R}-p_{n,R}$ have the $C^0$ decay in
\eqref{smooth:approximation-error} and the $C^r$ growth in
\eqref{smooth:approximant-growth}.  Interior interpolation on the fixed
pair $K\Subset K^+$ gives
\[
 \|p_{n+1,R}-p_{n,R}\|_{C^j(K)}\leq C e^{-\lambda n}.
\]
Thus the sequence converges in $C^j(K)$ to its already identified
uniform limit $q$.  Since $j$ is arbitrary, $q$ is smooth.  There is no
requirement that a single cutoff work for all $j$.  If a single sequence
converging in the $C^\infty$ topology is desired, choose cutoffs $R_k$
which work through order $k$, and choose $n_k$ so large that
$\|p_{n_k,R_k}-q\|_{C^k(K)}\leq2^{-k}$.  This diagonal sequence has that
property and still evaluates only finite polynomials.
\end{proof}

\subsection{Nondegeneracy and the strong bundles}

Smoothness of a leaf label does not by itself make its differential
surjective.  We use bounded holonomy Jacobians for this final step.

\begin{lemma}[Volume comparison for weak unstable holonomy]
\label{smooth:holonomy-volume}
Between two smooth $m$-dimensional local transversals to the weak
unstable foliation, its holonomy $h$ satisfies, locally,
\[
 c\,\Vol(E)\leq\Vol(h(E))\leq C\,\Vol(E)
\]
for measurable sets $E$, with $0<c\leq C<\infty$.
\end{lemma}

\begin{proof}
Strong unstable holonomy between smooth transversals is absolutely
continuous with locally bounded positive Jacobian; applying the same
statement to inverse holonomy supplies the reciprocal bound.  This is
the bounded-Jacobian theorem for horospheric foliations
\cite[Appendix, Theorem~5.1 and Definition~3.3]{Ballmann1995}.
We explain the passage to the weak foliation.

Thicken each transversal for a short interval by the flow, using a fixed
time translate of the target if necessary.  The resulting smooth
$(m+1)$-dimensional submanifolds are transversals to the strong unstable
foliation.  On suitably small flow boxes, strong holonomy has the form
\[
 H(a,t)=(h(a),t+\tau(a)).
\]
This follows from flow invariance and uniqueness of the crossing;
$\tau$ need only be continuous.  Choose a fixed short interval $I$ so
that $H$ is defined on $E\times I$ for every $E$ in a smaller source
patch.  In product coordinates, its image has the interval
$I+\tau(h^{-1}(y))$ above each $y\in h(E)$.  Fubini's theorem gives
product volume $|I|\Vol(h(E))$.  The smooth volume densities of both
flow-box parametrizations are bounded above and below.  The strong
holonomy comparison applied to $E\times I$ proves the displayed weak
comparison.  A finite chain of boxes proves the same assertion along
any compact holonomy path.
\end{proof}

\begin{theorem}[Smoothness of the Anosov splitting]
\label{smooth:splitting}
In the finite alternative of Theorem~\ref{alg:dichotomy}, the actual
bundles $E^u$ and $E^s$ are smooth on $SM$.
\end{theorem}

\begin{proof}
Fix a base point $x$ on the universal cover.  On its direction sphere,
the restriction $q|_{S_x\widetilde M}$ is locally the weak unstable
holonomy to $A$, followed by the smooth coordinate map $q_A$.  Both are
genuine $m$-dimensional transversals: a weak unstable tangent vector
with zero spatial projection is zero, and $TA=E^s$ is complementary to
$E^u\oplus\R X$.  The restriction is a local
homeomorphism, by the endpoint parametrizations.  By
Proposition~\ref{smooth:propagation} it is smooth, and by
Lemma~\ref{smooth:holonomy-volume} it satisfies a positive lower volume
comparison on every smaller patch.

Its derivative cannot be singular. Use smooth coordinate charts on
the source sphere and the target, with volume densities bounded above
and below. If the derivative at a point had rank less than $m$, choose
a unit target direction perpendicular to its image. Differentiability
then places the image of a coordinate ball of radius $r$ in a box with
widths $O(r)$ in the other $m-1$ directions and $o(r)$ in that direction.
Its volume is $o(r^m)$, contradicting the fixed positive lower volume
comparison for these balls. Thus $dq$ is surjective at every point of
$\mathcal U$.  The smooth bundle $\ker dq$ has dimension $m+1$ and
contains the tangent spaces of the actual weak unstable leaves, since
$q$ is constant on each such smooth leaf.  The dimensions agree, so
\[
 \ker dq=E^u\oplus\R X
 \quad\hbox{on }\mathcal U.
\]
Intersecting with the smooth contact distribution gives the smooth
strong unstable bundle
\[
 E^u=(\ker dq)\cap\ker\alpha.
\]

The coordinate patches of the full typical strong stable axis cover
all past endpoints except $\eta$.  A second typical axis with a distinct
future endpoint covers the omitted endpoint.  Such an axis exists
because typical points have full smooth measure, whereas a single
fixed-future weak leaf has positive codimension.  The preceding
argument therefore gives smoothness of $E^u$ everywhere on the cover.
The flip $v\mapsto-v$ exchanges the strong bundles, and hence gives
smoothness of $E^s$.  These are the actual invariant bundles, so they
are deck invariant and descend smoothly to $SM$.
\end{proof}
